% Standalone problem document, generated from the adjacent README.md.
% Compile with XeLaTeX. Mathematical content is maintained in Markdown.
\documentclass[11pt,a4paper]{article}
\usepackage[fontsize=10.5pt]{scrextend}
\usepackage[margin=22mm,headheight=15pt,headsep=9mm,footskip=11mm]{geometry}
\usepackage{fontspec}
\setmainfont{texgyrepagella}[Extension=.otf,UprightFont=*-regular,BoldFont=*-bold,ItalicFont=*-italic,BoldItalicFont=*-bolditalic]
\setsansfont{texgyreheros}[Extension=.otf,UprightFont=*-regular,BoldFont=*-bold,ItalicFont=*-italic,BoldItalicFont=*-bolditalic]
\setmonofont{texgyrecursor}[Extension=.otf,UprightFont=*-regular,BoldFont=*-bold,ItalicFont=*-italic,BoldItalicFont=*-bolditalic,Scale=MatchLowercase]
\usepackage{amsmath,amssymb}
\usepackage{unicode-math}
\setmathfont{texgyrepagella-math.otf}
\usepackage{xcolor,microtype,enumitem,needspace,fancyhdr}
\usepackage{bookmark}
\definecolor{ink}{HTML}{17344B}
\definecolor{accent}{HTML}{197D80}
\definecolor{muted}{HTML}{596773}
\hypersetup{colorlinks=true,linkcolor=accent,urlcolor=accent,pdftitle={MF-25 - Maximum
degree of real polynomials with a multiplication
budget},pdfauthor={Open Problems in Numerical Linear Algebra}}
\urlstyle{same}
\setlength{\parindent}{0pt}
\setlength{\parskip}{5pt plus 1pt minus 1pt}
\linespread{1.04}
\setlength{\emergencystretch}{3em}
\widowpenalty=10000
\clubpenalty=10000
\displaywidowpenalty=10000
\allowdisplaybreaks[1]
\setlist{nosep,leftmargin=1.5em}
\providecommand{\tightlist}{\setlength{\itemsep}{0pt}\setlength{\parskip}{0pt}}
\providecommand{\pandocbounded}[1]{#1}
\pagestyle{fancy}
\fancyhf{}
\fancyhead[L]{\sffamily\footnotesize\color{muted}OPEN PROBLEMS IN NUMERICAL LINEAR ALGEBRA}
\fancyhead[R]{\sffamily\bfseries\footnotesize\color{accent}MF-25}
\fancyfoot[L]{\parbox[b]{0.9\textwidth}{\sffamily\fontsize{8}{10}\selectfont\color{muted}Editorial ratings; a bounded search cannot certify that a problem remains unsolved.\\Literature check: 2026-11-27}}
\fancyfoot[R]{\sffamily\footnotesize\color{muted}\thepage}
\renewcommand{\headrulewidth}{0.3pt}
\renewcommand{\headrule}{\hbox to\headwidth{\color{accent}\leaders\hrule height \headrulewidth\hfill}}
\makeatletter
\renewcommand{\section}{\Needspace{5\baselineskip}\@startsection{section}{1}{0pt}{12pt plus 2pt minus 2pt}{4pt}{\sffamily\bfseries\large\color{ink}}}
\renewcommand{\subsection}{\Needspace{4\baselineskip}\@startsection{subsection}{2}{0pt}{9pt plus 1pt minus 1pt}{3pt}{\sffamily\bfseries\normalsize\color{ink}}}
\makeatother
\setcounter{secnumdepth}{0}
\begin{document}
{\sffamily\small\color{muted}Matrix functions and stability\par}
\vspace{3pt}
{\raggedright\hyphenpenalty=10000\exhyphenpenalty=10000\sffamily\bfseries\fontsize{21}{24}\selectfont\color{ink}Maximum
degree of real polynomials with a multiplication budget\par}
\vspace{7pt}
{\sffamily\small\textbf{Difficulty:} challenging\quad\textbf{Importance:} interesting
to the community\par}
{\sffamily\small\bfseries\color{accent}Status: Open\par}
\vspace{4pt}
{\color{accent}\hrule height 0.8pt}
\vspace{6pt}
\textbf{Rating rationale:} Real coefficient-space density and sharp
degree exclusions must both be established for arbitrary budgets. The
result would give a benchmark for designing low-cost matrix-function
approximations.

\hypertarget{problem-statement}{%
\subsection{Problem statement}\label{problem-statement}}

For a nonnegative integer \(r\), let \(\Pi_r\) be the set of real
polynomials of degree exactly \(r\). For a nonnegative integer \(m\),
let \(\Pi_{2^m}^*\) be the real polynomials computable from \(1,x\) by a
straight-line program with at most \(m\) nonscalar multiplications. Each
multiplication forms the product of two previously available
polynomials, while arbitrary real linear combinations are free. This is
the same arithmetic model as in
\href{https://github.com/ajt60gaibb/OpenProblemsInNLA/blob/main/matrix-functions-and-stability/MF-01/README.md}{MF-01}
and
\href{https://github.com/ajt60gaibb/OpenProblemsInNLA/blob/main/matrix-functions-and-stability/MF-14/README.md}{MF-14},
except that MF-14 uses complex coefficients and fixes the multiplication
budget to seven. Substituting a square matrix \(X\) for \(x\) gives an
evaluation scheme valid for matrices of arbitrary size.

Identify each polynomial of degree at most \(r\) with its monomial
coefficient vector in \(\mathbb R^{r+1}\). The bar denotes Euclidean
closure in this coefficient space. Define

\[d(m)=\max\left\{
r\in\mathbb N:\Pi_r\subseteq
\overline{\Pi_{2^m}^*\cap\Pi_r}
\right\}.\]

The inclusion at degree \(r\) means that for every \(p\in\Pi_r\) and
every \(\varepsilon>0\) there exists a \(q\in\Pi_{2^m}^*\cap\Pi_r\)
whose coefficient vector is within Euclidean distance \(\varepsilon\) of
that of \(p\). Thus every real polynomial of degree \(d(m)\) can be
approximated arbitrarily closely in its coefficients by polynomials of
the same degree that admit evaluation with at most \(m\) nonscalar
multiplications.

The classical Paterson--Stockmeyer algorithm provides a constructive
lower bound, while the dimension formula \(\dim(\Pi_{2^m}^*)=m^2\) for
\(m>2\) (Jarlebring and Lorentzon, Theorem 4.2) provides an upper bound.
Together, these give

\[\left\lfloor\frac{(m+2)^2}{4}\right\rfloor
\leq d(m)\leq m^2-1,\qquad m>2.\]

The lower bound follows from exact polynomial evaluation using the
Paterson--Stockmeyer scheme; the upper bound follows since the space of
polynomials of degree exactly \(r\) has dimension \(r+1\). The general
problem is to determine \(d(m)\), and we propose the following.

\textbf{Conjecture 1.} For every integer \(m\geq 4\),

\[d(m)=m^2-m.\]

Even if Conjecture 1 does not hold, the following weaker asymptotic
statement may still be true.

\textbf{Conjecture 2.}

\[\lim_{m\to\infty}\frac{d(m)}{m^2}=1.\]

\hypertarget{why-it-matters}{%
\subsection{Why it matters}\label{why-it-matters}}

Matrix products typically dominate the cost of evaluating a polynomial
of a large dense matrix. The conjecture would identify the greatest
degree at which any real polynomial approximant, for example one used to
approximate the matrix exponential, can be perturbed arbitrarily little
in its coefficients and evaluated within a prescribed product budget.

\hypertarget{related-problems}{%
\subsection{Related problems}\label{related-problems}}

\href{https://github.com/ajt60gaibb/OpenProblemsInNLA/blob/main/matrix-functions-and-stability/MF-01/README.md}{MF-01}
studies optimal approximation of the sign function under a fixed
multiplication budget. In contrast, the present problem concerns the
maximum polynomial degree achievable with a given number of
multiplications.

\href{https://github.com/ajt60gaibb/OpenProblemsInNLA/blob/main/matrix-functions-and-stability/MF-14/README.md}{MF-14}
concerns polynomial evaluation with seven matrix multiplications. Webb
(2026) established that the maximum degree in the complex, full-space
Zariski-closure model is 47. Since real Euclidean density within a
fixed-degree polynomial space implies complex Zariski density in the
full coefficient space, this result provides the upper bound

\[d(7)\leq47.\]

This improves the general dimension bound \(d(7)\leq48\), but does not
establish the conjectured equality \(d(7)=42\). The sense of
approximability the current conjecture leads to is different from that
of
\href{https://github.com/ajt60gaibb/OpenProblemsInNLA/blob/main/matrix-functions-and-stability/MF-14/README.md}{MF-14},
since higher order coefficients are here enforced to zero; which is may
improve practical usefulness.

See also
\href{https://github.com/ajt60gaibb/OpenProblemsInNLA/blob/main/references/webb-mf14-degree47-2026-09-17/proof.pdf}{Webb's
degree-47 proof}.

\hypertarget{references-and-status-evidence}{%
\subsection{References and status
evidence}\label{references-and-status-evidence}}

E. Jarlebring and G. Lorentzon, \emph{The Polynomial Set Associated with
a Fixed Number of Matrix-Matrix Multiplications}, SIAM Journal on
Applied Algebra and Geometry 10(3), 549--581 (2026).
\href{https://doi.org/10.1137/25M1779036}{DOI: 10.1137/25M1779036}.

M. S. Paterson and L. J. Stockmeyer, \emph{On the Number of Nonscalar
Multiplications Necessary to Evaluate Polynomials}, SIAM Journal on
Computing 2(1), 60--66 (1973).
\href{https://doi.org/10.1137/0202007}{DOI: 10.1137/0202007}.

\textbf{Rigorous results.} Theorem 4.2 of Jarlebring and Lorentzon
proves the dimension formula \(\dim(\Pi_{2^m}^*)=m^2\) for \(m>2\).
Section 5.1 gives an exact four-product construction with real
coefficients for every real polynomial of degree 12 and identifies 12 as
the optimal degree for coverage. In particular, the exact construction
gives the density inclusion at degree 12 in the formulation above.

\textbf{Numerical and empirical evidence.} The following results support
the conjecture:

\begin{itemize}
\tightlist
\item
  \textbf{\(d(5)=20\):} Real evaluation schemes for several degree-20
  polynomials were obtained using homotopy continuation, together with a
  combinatorial analysis of zero patterns in the evaluation schemes
  (Jarlebring and Lorentzon, §5, Conjecture 5.1). These findings are
  consistent with the earlier work of Sastre et al.~(2025, §4), who
  developed a five-multiplication evaluation scheme for degree-20
  polynomials, where realness of coefficients were not studied.
\item
  \textbf{\(d(6)=30\):} A real evaluation scheme for a degree-30
  polynomial was obtained using an iteratively regularized Newton
  method, together with a combinatorial analysis of zero patterns in the
  evaluation schemes (Jarlebring and Lorentzon, §5, Conjecture 5.2).
\end{itemize}

These results provide numerical support for the conjectured degrees but
do not establish Euclidean density or optimality. The numerical
experiments and evaluation schemes are available in the
\href{https://github.com/GustafLorentzon/polynomial-set-paper}{accompanying
software repository}.

\textbf{Relation to the published conjectures.} The conjecture above
generalizes the small-budget observations in Jarlebring and Lorentzon,
which states separate conjectures for five, six, and seven products. It
also uses a different notion of closure: Section 2.2 specifies Zariski
closure, and (5.1) takes the closure of the full output set before
testing coverage at a given degree. Here the approximating outputs must
have exactly the target degree and converge in real Euclidean
coefficient space.

\hypertarget{audit-2026-11-27}{%
\subsection{Audit --- 2026-11-27}\label{audit-2026-11-27}}

Checked the published SIAGA article's model and closure convention,
Theorem 4.2, the constructions and conjectures in §5, and the concluding
discussion. Also checked the local MF-14 resolution, §4 of Sastre et
al.'s degree-20 construction, and §6 of Alonso et al.'s preprint on
stable evaluation. A targeted search by the SIAGA title and by the
five-, six-, and seven-product degree claims located related
constructions and the complex MF-14 result.

This bounded check located no proof or counterexample for the stated
all-budget real, same-degree Euclidean conjecture.

\href{https://wseas.com/journals/mathematics/2025/b385106-036\%282025\%29.pdf}{Sastre
et al., §4} · \href{https://arxiv.org/html/2603.23143\#S6}{Alonso et
al., §6}.
\end{document}
